\section{Introduction}\label{sec:intro}

A deployed classifier often meets data that differ from its training data, and two kinds of change can happen at once.
The images may look different, which is known as \emph{domain shift}.
The new domain may also contain classes the model has never seen, which we call \emph{target-private classes}.
For example, an object recogniser trained on city scenes may be deployed in rural areas, where both the scenes and the objects differ~\citep{kundu2020cida}.

Existing methods handle only one of these changes.
Unsupervised domain adaptation (DA) aligns the two domains, but it assumes that both share one label set~\citep{ganin2016dann}.
Open-set and universal DA methods do detect new classes, but they merge all of them into a single ``unknown'' class~\citep{saito2018osbp,liu2019sta,you2019uan}.
Class-incremental (CI) learning adds new classes over time, yet it ignores domain shift and usually needs labelled data or the source data~\citep{li2017lwf,rebuffi2017icarl,castro2018e2e,dhar2019lwm}.

Kundu et~al.~\citep{kundu2020cida} define the CIDA setting to cover both changes under three constraints.
The first is that the source data are not available during adaptation, which makes the method \emph{source-free}.
The second is that the target images are unlabelled.
The third is that each target-private class has exactly one labelled image, which the paper calls the \emph{one-shot} setting.
Under these constraints, the goal is to classify both the shared and the new classes in the target domain.

Their method works in two stages.
The first stage trains the source model so that each class forms a compact Gaussian cluster in feature space.
It also trains the classifier on synthetic negative samples, which reserves room for classes the model has not yet seen.
The second stage adapts the model in a new feature space for the target, where learnable class centres called \emph{guides} provide the pseudo-labels.

We chose this paper under Option~1 of the assignment, which asks for the reimplementation of a paper without public code.
CIDA was published at ECCV 2020, and its authors did not release their code.
Our work makes four contributions.
\begin{itemize}
  \item \textbf{An implementation from scratch.} It has approximately 940 lines of Python for the method and 410 for analysis tools, not counting blank and comment lines. It also comes with 23 tests and job scripts for a SLURM cluster.
  \item \textbf{A reproduction of the method as written.} A configuration preset switches off all of our changes, so we can run the method exactly as the paper describes it. We use this reproduction to locate where accuracy is lost.
  \item \textbf{Four changes to the method.} Each change addresses one cause of the accuracy loss (\cref{sec:changes}).
  \item \textbf{Checks on unclear parts of the paper.} We list the settings the paper leaves open and three inconsistencies in it. The most important is a $3\sigma$ rule for negative samples that fails in high dimensions (\cref{sec:inconsistencies}).
\end{itemize}
